\documentclass[11pt,a4paper]{report}

\usepackage[margin=1in]{geometry}
\usepackage{amsmath,amssymb,amsfonts}
\usepackage{booktabs}
\usepackage{hyperref}
\usepackage{enumitem}
\usepackage{graphicx}
\usepackage{longtable}
\usepackage{bbm}

\hypersetup{
    colorlinks=true,
    linkcolor=blue,
    citecolor=blue,
    urlcolor=blue
}

\title{OD Estimation on the Sioux Falls Network\\[0.4em]
\large Methodology, Data Pipelines, and Machine-Learning Benchmarking}
\author{Implementation: \texttt{src/data/}, \texttt{src/ml/}, \texttt{generate\_od\_dataset.py}, \texttt{generate\_turning\_counts.py}, \texttt{train\_od\_estimator.py}}
\date{\today}

% --- Notation ---
\newcommand{\Tbase}{\mathbf{T}^{\mathrm{base}}}
\newcommand{\Tsyn}{\mathbf{T}^{\mathrm{syn}}}
\newcommand{\Tsynk}[1]{\mathbf{T}^{\mathrm{syn},(#1)}}
\newcommand{\Tsycand}{\mathbf{T}^{\mathrm{syn,cand}}}
\newcommand{\Xipf}[1]{\mathbf{X}^{(#1)}}
\newcommand{\mathbfp}{\mathbf{p}}
\newcommand{\mathbfa}{\mathbf{a}}
\newcommand{\mathbfS}{\mathbf{S}}
\newcommand{\mathbfM}{\mathbf{M}}
\newcommand{\mathbfy}{\mathbf{y}}
\newcommand{\mathbfx}{\mathbf{x}}
\newcommand{\mathbfX}{\mathbf{X}}
\newcommand{\mathbfY}{\mathbf{Y}}
\newcommand{\Aturn}{\mathbf{A}_{\mathrm{turn}}}
\newcommand{\mathbfT}{\mathbf{T}}
\newcommand{\mathbfA}{\mathbf{A}}
\newcommand{\mathbfw}{\mathbf{w}}
\newcommand{\mathbfq}{\mathbf{q}}
\newcommand{\mathbfz}{\mathbf{z}}

\begin{document}
\maketitle

\begin{abstract}
This document describes the full pipeline for origin--destination (OD) estimation on the Sioux Falls network, organized in three chapters: (i)~synthetic OD matrix generation and evaluation, (ii)~probabilistic turning-count observation via a route-choice-aware linear assignment operator, and (iii)~machine-learning model development and benchmarking for the inverse problem. Notation is fixed throughout. The synthetic generator builds a distance-dependent gravity prior on full off-diagonal support, draws Dirichlet samples for compositional diversity, applies Gamma heavy-tailed reweighting for realistic flow skewness, scales demand, and balances with iterative proportional fitting (IPF). Candidate matrices are then accepted or rejected by dataset-level quality filters (sparsity, cell-wise entropy, pairwise Frobenius diversity, correlation bounds, and trip-length consistency). Turning counts are generated by $K$-shortest-path enumeration, multinomial logit route choice, and a fractional assignment matrix $\Aturn$, with optional Poisson observation noise. Implementations live under \texttt{src/data/} and \texttt{src/ml/}. The design supports datasets of up to $10^6$ samples using batched, memory-mapped I/O.
\end{abstract}

\tableofcontents
\newpage

%=============================================================================
\chapter{Synthetic Data Generation and Evaluation}
\label{ch:synthetic}
%=============================================================================

\section{Problem setting and pipeline overview}
\label{sec:ch1-problem}

\subsection{The OD matrix}

Let $N=24$ be the number of traffic analysis zones in Sioux Falls. The \emph{base} OD matrix
\begin{equation}
    \Tbase = \bigl(T^{\mathrm{base}}_{ij}\bigr)_{i,j=1}^{N} \in \mathbb{R}_{\ge 0}^{N \times N}
\end{equation}
records daily trips from origin zone $i$ to destination zone $j$. We observe a single matrix \texttt{EstimatedODMatrix.npy}; the goal is to train models that recover OD demand from turning-count observations.

\paragraph{Marginals and grand total.}
\begin{align}
    p_i^{\mathrm{base}} &= \sum_{j=1}^{N} T^{\mathrm{base}}_{ij}
    && \text{(productions)}, \\
    a_j^{\mathrm{base}} &= \sum_{i=1}^{N} T^{\mathrm{base}}_{ij}
    && \text{(attractions)}, \\
    G &= \sum_{i,j} T^{\mathrm{base}}_{ij}
    = \sum_i p_i^{\mathrm{base}} = \sum_j a_j^{\mathrm{base}}.
\end{align}

\subsection{Why synthetic data?}

Machine learning requires many labeled examples. We generate $K_{\mathrm{samples}}$ synthetic matrices $\{\Tsynk{k}\}_{k=1}^{K_{\mathrm{samples}}}$ that:
\begin{itemize}[noitemsep]
    \item preserve total demand $G$ and forbid intrazonal trips ($T_{ii}=0$);
    \item admit full off-diagonal support $i\neq j$, guided by a gravity-style prior;
    \item exhibit heavy-tailed cell flows (many tiny OD pairs, few large ones);
    \item pass dataset-level diversity and realism filters before inclusion;
    \item pair with turning counts via a known forward model for supervised training.
\end{itemize}

\subsection{End-to-end pipeline}

\begin{figure}[ht]
\centering
\fbox{\begin{minipage}{0.94\linewidth}
\small
\textbf{Chapter 1:}
$\Tbase \to$ prior $(T+\varepsilon) \to$ Dirichlet $\to$ Gamma weighting $\to$ scale $\to$ IPF
$\to$ quality filters $\to$ $\{\Tsynk{k}\}$\\[0.3em]
\textbf{Chapter 2:} $\Tsynk{k} \to \mathbfy^{(k)} \to \mathbfx^{(k)} = \Aturn\mathbfy^{(k)} + \boldsymbol{\epsilon}^{(k)}$ \quad ($K$-paths + logit + fractional $\Aturn$)\\[0.3em]
\textbf{Chapter 3:} $\mathbfX^{\mathrm{obs}} \to f(\cdot) \to \hat{\mathbfY}$ \quad (information analysis $\to$ baselines $\to$ ML benchmark)
\end{minipage}}
\caption{Three-chapter pipeline from base OD to ML-ready $(\mathbfX,\mathbfY)$ pairs.}
\label{fig:pipeline_overview}
\end{figure}

\subsection{Notation summary}

\begin{table}[ht]
\centering
\caption{Primary notation used throughout this document.}
\label{tab:notation}
\small
\begin{tabular}{@{}ll@{}}
\toprule
\textbf{Symbol} & \textbf{Meaning} \\
\midrule
$\Tbase$ & Observed base OD matrix \\
$\Tsyn$, $\Tsynk{k}$ & Synthetic OD (single / $k$-th replicate) \\
$\mathbfp^{\mathrm{base}}$, $\mathbfa^{\mathrm{base}}$ & Base productions and attractions \\
$\mathbfp^{\mathrm{tgt}}$, $\mathbfa^{\mathrm{tgt}}$ & IPF target marginals (sum to $G$) \\
$p_i^{\mathrm{syn}}$, $a_j^{\mathrm{syn}}$ & Achieved marginals of $\Tsyn$ \\
$\mathbfp^{\mathrm{prior}}$, $\varepsilon_{ij}$ & Gravity prior probabilities / distance-dependent additive mass \\
$d_{ij}$, $\beta$, $\lambda$ & Inter-zone distance; prior scale and decay length \\
$\tilde{\mathbfp}^{\mathrm{Dir}}$, $\mathbfw$ & Dirichlet sample; Gamma weights \\
$\mathbfp^{\mathrm{prob}}$ & Sampling prior on support $|\mathcal{S}|=N(N-1)=552$ \\
$\mathbfS$, $\Xipf{t}$ & IPF seed matrix / iterate at step $t$ \\
$\mathbfM$ & Off-diagonal support mask ($M_{ij}=\mathbbm{1}[i\neq j]$) \\
$\mathbfy$, $\mathbfY$ & Flattened OD vector / batch ($N^2=576$) \\
$\mathbfx$, $\mathbfX$ & Turning-count vector / batch ($M=178$); clean or noisy \\
$\mathbfx^{\mathrm{clean}}$, $\mathbfx^{\mathrm{obs}}$ & Clean turning counts / noisy observations \\
$\Aturn$ & Fractional turning assignment matrix $(178 \times 576)$ \\
$\Aturn^{\mathrm{bin}}$ & Binary (shortest-path) assignment matrix \\
$\mathcal{R}_{ij}$ & Set of enumerated routes for OD pair $(i,j)$ \\
$K$ & Number of shortest simple paths per OD pair \\
$\theta$ & Logit sensitivity (cost scale in route choice) \\
$c_r$, $U_r$, $P_r$ & Route cost, utility, choice probability \\
$\mathbf{b}_r$ & Binary turn-incidence vector for route $r$ \\
$\alpha$, $\delta$ & Dirichlet concentration / marginal perturbation level \\
$a_{\gamma}$, $b_{\gamma}$ & Gamma shape and scale for heavy-tailed weighting \\
$H_{ij}$ & Cell-wise Shannon entropy across the dataset \\
\bottomrule
\end{tabular}
\end{table}

\subsection{Codebase organization}

Separating \textbf{data generation} from \textbf{model estimation} is advisable: the two stages have different dependencies (NetworkX vs.\ PyTorch/XGBoost), different scaling bottlenecks, and different release cadences.

\begin{table}[ht]
\centering
\caption{Recommended \texttt{src/} package layout.}
\label{tab:code_layout}
\small
\begin{tabular}{@{}p{3.2cm}p{5.2cm}p{4.8cm}@{}}
\toprule
\textbf{Package} & \textbf{Responsibility} & \textbf{Entry scripts} \\
\midrule
\texttt{src/data/} &
Synthetic OD, turning counts, $\Aturn$ &
\texttt{generate\_od\_dataset.py}, \texttt{generate\_turning\_counts.py} \\
\texttt{src/ml/} &
Information analysis, training, benchmarking &
\texttt{train\_od\_estimator.py} \\
\texttt{src/\_\_init\_\_.py} &
Shared constants ($N=24$, paths) &
--- \\
\bottomrule
\end{tabular}
\end{table}

\paragraph{Why separate packages.}
\begin{itemize}[noitemsep]
    \item \textbf{Dependency isolation:} data generation requires no GPU/PyTorch; ML stage optionally does.
    \item \textbf{Memory profile:} Chapter~2 uses memory-mapped batch I/O; Chapter~3 uses train/val/test arrays.
    \item \textbf{Testing:} data-pipeline tests and ML tests run independently.
    \item \textbf{Clarity:} inverse-problem code does not import route enumeration modules.
\end{itemize}


\section{Structural constraints and design rationale}
\label{sec:ch1-constraints}

This chapter treats OD synthesis as a constrained random measure on zone pairs. Four modelling decisions shape the generator: (i)~which cells may carry flow, (ii)~how a prior encodes spatial structure, (iii)~how samples explore the probability simplex, and (iv)~how dataset-level filters enforce diversity and realism.

\subsection{Support and hard constraints}

Unlike an earlier design that restricted sampling to cells with positive base flow, the current support is \emph{full off-diagonal}:
\begin{equation}
    \mathcal{S} = \{(i,j): i\neq j\}, \qquad
    |\mathcal{S}| = N(N-1) = 552, \qquad
    M_{ij} = \mathbbm{1}[i\neq j].
    \label{eq:full_support}
\end{equation}
Hard constraints enforced for every synthetic replicate are
\begin{align}
    T_{ii}^{\mathrm{syn}} &= 0 && \text{(no intrazonal trips)}, \\
    \sum_{i,j} T_{ij}^{\mathrm{syn}} &= G && \text{(conserved demand)}, \\
    T_{ij}^{\mathrm{syn}} &\ge 0 && \text{(nonnegative flows)}.
\end{align}
Allowing $i\neq j$ cells that were zero in $\Tbase$ is deliberate: ML estimators of underdetermined OD problems must not be trained under an artificially narrow support that coincides with one observed matrix. The gravity prior of Section~\ref{sec:ch1-prior} then softly concentrates mass on short, historically active pairs without forbidding exploration of longer links.

\begin{table}[ht]
\centering
\caption{Generative approaches compared (Sioux Falls). The production pipeline is prior + Dirichlet + Gamma + IPF + quality filters.}
\label{tab:gen_compare}
\small
\begin{tabular}{@{}p{2.6cm}p{3.5cm}p{3.5cm}p{3.5cm}@{}}
\toprule
\textbf{Criterion} &
\textbf{Legacy noise} &
\textbf{Dirichlet+IPF only} &
\textbf{Current (full)} \\
\midrule
Support &
$\{T^{\mathrm{base}}>0\}$ &
$\{T^{\mathrm{base}}>0\}$ &
Full $i\neq j$ + distance prior \\
Composition &
$\eta_{ij}\sim\mathcal{U}(0.8,1.2)$ &
$\mathrm{Dir}(\alpha\mathbfp)$ &
Dirichlet then Gamma reweight \\
Skewness &
Mild &
Mild--moderate &
Heavy-tailed (many tiny / few large) \\
Diversity control &
Weak &
$\alpha$ only &
$\alpha$, $(a_\gamma,b_\gamma)$, filters \\
Typical $\rho$ vs.\ base &
$\approx 0.97$ &
$\approx 0.74$ ($\alpha=500$) &
Bounded by filter band $[\rho_{\min},\rho_{\max}]$ \\
\bottomrule
\end{tabular}
\end{table}

\begin{table}[ht]
\centering
\caption{Pros and cons of support choices.}
\label{tab:pros_support}
\small
\begin{tabular}{@{}p{3.4cm}p{5.0cm}p{5.0cm}@{}}
\toprule
\textbf{Choice} & \textbf{Pros} & \textbf{Cons} \\
\midrule
Support $\{T^{\mathrm{base}}>0\}$ &
Matches observed sparsity; stable IPF &
Cannot invent missing OD pairs; ML may overfit support \\
Full off-diagonal $\mathcal{S}$ &
Richer label space; consistent with zoning &
Needs a prior; more near-zero cells \\
\bottomrule
\end{tabular}
\end{table}


\section{Generative methodology}
\label{sec:ch1-generation}

The production generator proceeds in six stages:
\begin{equation}
\begin{aligned}
\Tbase
&\xrightarrow{\text{prior}}
\mathbfp^{\mathrm{prior}}
\xrightarrow{\mathrm{Dir}}
\tilde{\mathbfp}^{\mathrm{Dir}}
\xrightarrow{\mathrm{Gamma}}
\mathbfp^{\mathrm{prob}}
\\[-0.2em]
&\xrightarrow{\text{scale}}
\mathbfS
\xrightarrow{\mathrm{IPF}}
\Tsycand
\xrightarrow{\text{filters}}
\{\Tsynk{k}\}.
\end{aligned}
\label{eq:pipeline_stages}
\end{equation}

\subsection{Stage 1: Base OD analysis}
\label{sec:ch1-base}

Given $\Tbase$, compute grand total $G=\sum_{i,j}T^{\mathrm{base}}_{ij}$, productions $\mathbfp^{\mathrm{base}}$ and attractions $\mathbfa^{\mathrm{base}}$, and descriptive statistics (sparsity, top OD pairs). Total demand $G$ is held fixed in subsequent sampling; only the \emph{composition} of trips is randomized. Implementation: \texttt{src/data/od\_analysis.py}.

\subsection{Stage 2: Distance-dependent gravity prior}
\label{sec:ch1-prior}

Let $d_{ij}$ be Euclidean distance between zone centroids (junction coordinates from the SUMO network). Define additive prior mass
\begin{equation}
    \varepsilon_{ij} =
    \begin{cases}
      \beta\,\mathrm{e}^{-d_{ij}/\lambda}, & i\neq j, \\
      0, & i=j,
    \end{cases}
    \label{eq:epsilon}
\end{equation}
and the prior trip table
\begin{equation}
    T^{\mathrm{prior}}_{ij} = T^{\mathrm{base}}_{ij} + \varepsilon_{ij}, \qquad
    p^{\mathrm{prior}}_{ij}
    = \frac{T^{\mathrm{prior}}_{ij}}{\sum_{k,\ell} T^{\mathrm{prior}}_{k\ell}}.
    \label{eq:prior}
\end{equation}
Defaults: $\beta=1.0$, $\lambda=500\,\mathrm{m}$ (order of the median inter-zone distance on Sioux Falls). The exponential decay encodes a gravity-style preference for shorter trips while retaining the observed base structure in $T^{\mathrm{base}}$. Importantly, $G$ used for scaling remains $\sum T^{\mathrm{base}}$, not $\sum T^{\mathrm{prior}}$---the prior only shapes probabilities.

\paragraph{Why this prior.}
Base-only normalization $T^{\mathrm{base}}/G$ leaves historically inactive OD pairs at probability zero, so Dirichlet never places mass there. Adding $\varepsilon_{ij}>0$ opens full off-diagonal support in a continuous, distance-aware manner without imposing a hard cut-off.

\begin{table}[ht]
\centering
\caption{Pros and cons of prior constructions.}
\label{tab:pros_prior}
\small
\begin{tabular}{@{}p{3.2cm}p{5.1cm}p{5.1cm}@{}}
\toprule
\textbf{Method} & \textbf{Pros} & \textbf{Cons} \\
\midrule
$T^{\mathrm{base}}/G$ only &
Faithful to survey; simple &
Zeros remain forever; no gravity signal \\
Uniform off-diagonal &
Maximal exploration &
Ignores geography and base hierarchy \\
$T+\varepsilon$ with \eqref{eq:epsilon} &
Smooth full support; short trips preferred &
Requires $(\beta,\lambda)$; mild leakage onto long OD pairs \\
\bottomrule
\end{tabular}
\end{table}

\subsection{Stage 3: Dirichlet sampling on the simplex}
\label{sec:ch1-dirichlet}

Let $\mathbfp^{\mathrm{prior}}_{\mathcal{S}}\in\Delta^{|\mathcal{S}|-1}$ be the prior restricted to support cells. Draw
\begin{equation}
    \tilde{\mathbfp}^{\mathrm{Dir}}
    \sim \mathrm{Dirichlet}\!\bigl(\alpha\,\mathbfp^{\mathrm{prior}}_{\mathcal{S}}\bigr).
    \label{eq:dirichlet}
\end{equation}
The Dirichlet density on the probability simplex $\Delta^{d-1}$ with parameter $\boldsymbol{\alpha}=(\alpha_1,\ldots,\alpha_d)$, $\alpha_k>0$, is
\begin{equation}
    f(\mathbfp)
    = \frac{\Gamma\!\bigl(\sum_k \alpha_k\bigr)}{\prod_k \Gamma(\alpha_k)}
    \prod_{k=1}^{d} p_k^{\alpha_k-1},
    \qquad p_k>0,\; \sum_k p_k=1.
    \label{eq:dirichlet_pdf}
\end{equation}
With the proportional parameterization $\alpha_k = \alpha\, p^{\mathrm{prior}}_k$, the mean is
\begin{equation}
    \mathbb{E}[\tilde{p}_k^{\mathrm{Dir}}] = p^{\mathrm{prior}}_k,
\end{equation}
and the variance of each coordinate decreases as $\alpha$ increases:
\begin{equation}
    \mathrm{Var}(\tilde{p}_k^{\mathrm{Dir}})
    = \frac{p^{\mathrm{prior}}_k\bigl(1-p^{\mathrm{prior}}_k\bigr)}{\alpha+1}.
    \label{eq:dirichlet_var}
\end{equation}
Thus $\alpha$ is a single scalar knob trading off realism against compositional diversity.

\paragraph{Why Dirichlet.}
Trip tables are compositional: shares must sum to one under fixed $G$. The Dirichlet is the conjugate / natural distribution on that simplex; proportional means preserve the prior ranking of OD pairs in expectation while allowing controlled random reallocation of mass. Independent Gaussian noise on cells would violate the simplex (or require clumsy post-normalization) and ignore the compositional constraint.

\begin{table}[ht]
\centering
\caption{Dirichlet concentration $\alpha$: empirical trade-offs on Sioux Falls
($K_{\mathrm{samples}}=5000$, legacy support study; production uses the same qualitative regime).}
\label{tab:alpha}
\small
\begin{tabular}{@{}llll@{}}
\toprule
$\alpha$ & Mean $\rho$ vs.\ base & Pairwise diversity & Recommendation \\
\midrule
100  & lower ($\sim$0.65) & highest & exploratory / high-diversity ML \\
250  & moderate & high & alternative if $\alpha=500$ too conservative \\
500  & $\approx 0.74$ & $\approx 545$ & default concentration \\
1000 & higher & moderate & closer to base \\
2000 & highest ($\sim$0.90+) & lowest & ablation only \\
\bottomrule
\end{tabular}
\end{table}

\begin{table}[ht]
\centering
\caption{Pros and cons of compositional sampling methods.}
\label{tab:pros_dirichlet}
\small
\begin{tabular}{@{}p{3.2cm}p{5.1cm}p{5.1cm}@{}}
\toprule
\textbf{Method} & \textbf{Pros} & \textbf{Cons} \\
\midrule
Dirichlet \eqref{eq:dirichlet} &
Correct support (simplex); mean = prior; one knob $\alpha$ &
Does not alone create extreme flow skewness \\
Independent lognormal cells + normalize &
Easy heavy tails &
Normalization couples cells awkwardly; hard to set means \\
Additive Gaussian on $\Tbase$ &
Simple &
Can leave simplex/negativity; clones of base at high SNR \\
\bottomrule
\end{tabular}
\end{table}

\subsection{Stage 4: Gamma heavy-tailed reweighting}
\label{sec:ch1-gamma}

Empirical OD matrices are highly skewed: most zone pairs carry small flows, while a few corridors dominate. Dirichlet draws alone tend to be smoother than this. After Dirichlet, we therefore apply independent positive weights
\begin{equation}
    w_{ij} \sim \mathrm{Gamma}(a_{\gamma},\, b_{\gamma}),
    \qquad i\neq j,
    \label{eq:gamma_w}
\end{equation}
in the shape--scale parameterization with density
\begin{equation}
    f(w)
    = \frac{1}{\Gamma(a_{\gamma})\, b_{\gamma}^{a_{\gamma}}}
    w^{a_{\gamma}-1}\,\mathrm{e}^{-w/b_{\gamma}},
    \qquad w>0,
    \label{eq:gamma_pdf}
\end{equation}
mean $\mathbb{E}[w]=a_{\gamma}b_{\gamma}$ and variance $\mathrm{Var}(w)=a_{\gamma}b_{\gamma}^2$. Reweight and renormalize:
\begin{equation}
    p_{ij}^{\mathrm{prob}}
    = \frac{w_{ij}\,\tilde{p}_{ij}^{\mathrm{Dir}}}
           {\displaystyle\sum_{k\neq\ell} w_{k\ell}\,\tilde{p}_{k\ell}^{\mathrm{Dir}}},
    \qquad
    p_{ii}^{\mathrm{prob}}=0.
    \label{eq:gamma_reweight}
\end{equation}
Defaults: $a_{\gamma}=0.5$, $b_{\gamma}=1.0$. For $a_{\gamma}<1$, the Gamma density diverges as $w\to 0^+$ and retains a long right tail, yielding many tiny OD shares and occasional large weights---exactly the heavy-tailed pattern desired for synthetic tables. Disable with \texttt{--no-gamma-mask}.

\paragraph{Why Gamma (not LogNormal / Pareto alone).}
Gamma provides strictly positive weights with a single, interpretable shape controlling near-zero mass. LogNormal is also heavy-tailed but less flexible near zero for the same two-parameter budget; Pareto concentrates too aggressively on extremes and can destabilize IPF. Mixing Dirichlet (composition) with Gamma (magnitude skew) separates ``where mass can move'' from ``how uneven magnitudes are.''

Note that a Dirichlet distribution itself can be constructed by normalizing independent Gamma random variables. However, those Gamma variables determine only the compositional uncertainty implied by the prior. The additional Gamma reweighting introduced here is an independent multiplicative process that models empirical flow heterogeneity, producing the heavy-tailed corridor hierarchy observed in real transportation networks without altering the underlying prior expectation.

\begin{table}[ht]
\centering
\caption{Pros and cons of heavy-tailed reweighting.}
\label{tab:pros_gamma}
\small
\begin{tabular}{@{}p{3.0cm}p{5.2cm}p{5.2cm}@{}}
\toprule
\textbf{Method} & \textbf{Pros} & \textbf{Cons} \\
\midrule
None (Dirichlet only) &
Stable; fewer hyperparameters &
Over-smooth OD hierarchy \\
Gamma $a_{\gamma}<1$ &
Many tiny / few large flows; positive weights &
Extra $(a_{\gamma},b_{\gamma})$; can dilute base pattern \\
LogNormal &
Classic heavy tail &
Weaker near-zero density control \\
Pareto / power-law &
Extreme corridors &
Can break IPF / correlation filters \\
\bottomrule
\end{tabular}
\end{table}

\subsection{Stage 5: Demand scaling and optional marginal perturbation}

Form the IPF seed
\begin{equation}
    S_{ij} = G\, p_{ij}^{\mathrm{prob}}, \qquad S_{ii}=0.
\end{equation}
Optional multiplicative shocks on zone totals ($\delta\ge 0$; production often uses $\delta=0$ when filters already control diversity):
\begin{align}
    \tilde{p}_i &= p_i^{\mathrm{base}}(1+\xi_i^p), &
    \xi_i^p &\sim \mathcal{U}(-\delta,\delta), \\
    \tilde{a}_j &= a_j^{\mathrm{base}}(1+\xi_j^a), &
    \xi_j^a &\sim \mathcal{U}(-\delta,\delta), \\
    p_i^{\mathrm{tgt}} &= G\,\tilde{p}_i\Big/\textstyle\sum_{i'}\tilde{p}_{i'}, &
    a_j^{\mathrm{tgt}} &= G\,\tilde{a}_j\Big/\textstyle\sum_{j'}\tilde{a}_{j'}.
\end{align}
With $\delta=0$, targets equal the base marginals (rescaled to $G$).

\begin{table}[ht]
\centering
\caption{Pros and cons of marginal perturbation.}
\label{tab:pros_delta}
\small
\begin{tabular}{@{}p{3.2cm}p{5.1cm}p{5.1cm}@{}}
\toprule
\textbf{Setting} & \textbf{Pros} & \textbf{Cons} \\
\midrule
$\delta=0$ &
Stable productions/attractions; clear ablation &
Less zone-level diversity \\
$\delta>0$ (e.g.\ 0.20) &
Realistic zone growth/decline &
Can fight gravity prior / trip-length filters \\
\bottomrule
\end{tabular}
\end{table}

\subsection{Stage 6: Iterative proportional fitting (IPF)}

IPF (Furness / RAS) balances the seed to the target marginals. With $\Xipf{0}=\mathbfS$,
\begin{align}
    X^{(t+\frac{1}{2})}_{ij}
    &= X^{(t)}_{ij}\cdot
    \frac{p_i^{\mathrm{tgt}}}{\sum_{j'}X^{(t)}_{ij'}},
    \label{eq:ipf_row} \\
    X^{(t+1)}_{ij}
    &= X^{(t+\frac{1}{2})}_{ij}\cdot
    \frac{a_j^{\mathrm{tgt}}}{\sum_{i'}X^{(t+\frac{1}{2})}_{i'j}}.
    \label{eq:ipf_col}
\end{align}
Stop when marginal error $e^{(t)}<\tau=10^{-3}$ or after $t_{\max}=1000$ iterations; then $\Tsycand=\Xipf{t}\odot\mathbfM$ with zero diagonal.

\begin{table}[ht]
\centering
\caption{Pros and cons of balancing methods.}
\label{tab:pros_ipf}
\small
\begin{tabular}{@{}p{3.2cm}p{5.1cm}p{5.1cm}@{}}
\toprule
\textbf{Method} & \textbf{Pros} & \textbf{Cons} \\
\midrule
IPF / Furness &
Classic; fast; preserves zeros of the seed pattern up to scaling &
Requires feasible marginals; iterative \\
Entropy maximization &
Elegant dual interpretation &
Heavier numerically for $10^6$ samples \\
No balancing &
Cheap &
Inconsistent zone totals \\
\bottomrule
\end{tabular}
\end{table}

\begin{figure}[ht]
\centering
\fbox{\begin{minipage}{0.94\linewidth}
\small
\textbf{Algorithm 1:} Prior + Dirichlet + Gamma + IPF (one candidate)\\[0.4em]
\begin{enumerate}[leftmargin=*, nosep]
    \item Build $T^{\mathrm{prior}}=T^{\mathrm{base}}+\varepsilon$ with \eqref{eq:epsilon}--\eqref{eq:prior}; mask $\mathbfM$
    \item Draw $\tilde{\mathbfp}^{\mathrm{Dir}}\sim\mathrm{Dir}(\alpha\mathbfp^{\mathrm{prior}}_{\mathcal{S}})$
    \item Sample $w_{ij}\sim\mathrm{Gamma}(a_{\gamma},b_{\gamma})$; form $\mathbfp^{\mathrm{prob}}$ by \eqref{eq:gamma_reweight}
    \item Seed $\mathbfS=G\,\mathrm{reshape}(\mathbfp^{\mathrm{prob}})$; build targets $(\mathbfp^{\mathrm{tgt}},\mathbfa^{\mathrm{tgt}})$
    \item $\Tsycand\leftarrow\mathrm{IPF}(\mathbfS,\mathbfp^{\mathrm{tgt}},\mathbfa^{\mathrm{tgt}})$; project onto $\mathbfM$
\end{enumerate}
\end{minipage}}
\caption{Candidate generator (\texttt{src/data/od\_generator.py}, \texttt{dirichlet\_sampler.py}, \texttt{sparse\_mask.py}).}
\label{alg:dirichlet_ipf}
\end{figure}


\section{Quality filters (dataset acceptance)}
\label{sec:ch1-filters}

Candidate generation alone does not guarantee a useful \emph{dataset}: replicates can cluster, clone $\Tbase$, or violate trip-length realism. We therefore generate an oversized pool of size
\begin{equation}
    K_{\mathrm{pool}} = \lceil \kappa\, K_{\mathrm{samples}}\rceil
    \qquad (\text{default }\kappa=10)
\end{equation}
and accept a subset of size $K_{\mathrm{samples}}$ by sequential filters (implementation: \texttt{src/data/quality\_filters.py}).

\paragraph{Sparsity band.}
Let $s(\mathbfT)=\frac{1}{N^2}\sum_{i,j}\mathbbm{1}[T_{ij}\le 10^{-12}]$. Require
$s_{\min}\le s(\Tsycand)\le s_{\max}$ (defaults $0$ and $0.95$).

\paragraph{Correlation band.}
With $\rho=\mathrm{corr}(\mathrm{vec}(\Tbase),\mathrm{vec}(\Tsycand))$, require
$\rho_{\min}\le\rho\le\rho_{\max}$ (defaults $0.30$ and $0.95$) to reject both clones and wild outliers.

\paragraph{Trip-length consistency.}
Flow-weighted mean trip length
\begin{equation}
    \bar{\ell}(\mathbfT)
    = \frac{\sum_{i,j} T_{ij}\, d_{ij}}{\sum_{i,j} T_{ij}}
\end{equation}
must satisfy
$\bigl|\bar{\ell}(\Tsycand)-\bar{\ell}(\Tbase)\bigr|/\bar{\ell}(\Tbase)
\le \eta_{\ell}$ (default $\eta_{\ell}=0.35$).

\paragraph{Cell-wise entropy ranking (Constraint~3).}
For cell $(i,j)$ across a batch $\{\Tsynk{k}\}_{k=1}^{K}$, discretize flows into $B$ bins and compute Shannon entropy
\begin{equation}
    H_{ij}
    = -\sum_{b=1}^{B} \hat{\pi}_{ij,b}\,\log \hat{\pi}_{ij,b},
    \qquad
    \overline{H}=\frac{1}{|\mathcal{S}|}\sum_{(i,j)\in\mathcal{S}} H_{ij}.
    \label{eq:cell_entropy}
\end{equation}
Candidates are ranked to prefer increments in $\overline{H}$ (diverse marginal histograms per OD cell).

\paragraph{Pairwise Frobenius diversity (Constraint~4).}
Reject $\mathbfT$ if there exists an already accepted $\mathbfT'$ with
\begin{equation}
    \|\mathbfT-\mathbfT'\|_F
    < \tau_F,
    \qquad
    \|\mathbfA\|_F=\sqrt{\sum_{i,j} A_{ij}^2}.
    \label{eq:frobenius}
\end{equation}
If $\tau_F$ is unset, set $\tau_F = 0.35\times\mathrm{median}\{\|\mathbfT^{(a)}-\mathbfT^{(b)}\|_F\}$ over a sample of candidate pairs.

\begin{table}[ht]
\centering
\caption{Pros and cons of quality filters.}
\label{tab:pros_filters}
\small
\begin{tabular}{@{}p{3.2cm}p{5.1cm}p{5.1cm}@{}}
\toprule
\textbf{Filter} & \textbf{Pros} & \textbf{Cons} \\
\midrule
Sparsity / corr.\ bands &
Cheap realism gates &
Thresholds are network-specific \\
Trip-length check &
Geography-aware acceptance &
Needs distances; couples to prior \\
Cell entropy $H_{ij}$ &
Encourages rich per-cell variation &
Binning choices; soft ranking \\
Frobenius diversity &
Explicit anti-cloning of accept set &
$O(K_{\mathrm{acc}})$ checks; needs oversized pool \\
No filters &
Maximum throughput &
Redundant / unrealistic batches \\
\bottomrule
\end{tabular}
\end{table}


\section{Evaluation metrics and hyperparameter selection}
\label{sec:ch1-metrics}

\subsection{Per-matrix metrics}

For each $\Tsyn$, with $d_{ij}=T^{\mathrm{syn}}_{ij}-T^{\mathrm{base}}_{ij}$:
\begin{align}
    \mathrm{MAE} &= \frac{1}{N^2}\sum_{i,j}|d_{ij}|, &
    \mathrm{RMSE} &= \sqrt{\frac{1}{N^2}\sum_{i,j}d_{ij}^2}, \\
    \mathrm{RelErr} &= \frac{\sum_{i,j}|d_{ij}|}{G}, &
    \rho &= \mathrm{corr}(\mathrm{vec}(\Tbase),\mathrm{vec}(\Tsyn)).
\end{align}
Marginal drift: $E_p = \frac{1}{N}\sum_i|({p_i^{\mathrm{syn}}-p_i^{\mathrm{base}}})/({p_i^{\mathrm{base}}+10^{-9}})|$ (and analogously $E_a$). Structural validity requires zero diagonal mass and nonnegative flows on $\mathcal{S}$.

\subsection{Dataset-level metrics and hyperparameter search}

For batch $\{\Tsynk{k}\}$: pairwise diversity (L1 / Frobenius), per-OD mean/std/CV, demand consistency $\mathrm{CV}_G\approx 0$, mean cell entropy $\overline{H}$, PCA explained variance, and KMeans inertia. For $K_{\mathrm{samples}}>10^5$, metrics use a random subset (default 10{,}000).

\texttt{src/data/hyperparameter\_search.py} evaluates $(\alpha,\delta)$ pairs and ranks by
\begin{equation}
    \mathrm{score} = \frac{\bar{d}}{200} + \overline{\mathrm{CV}}_{\mathrm{cell}} + \frac{n_{90\%}}{20} - 2\max(0,\bar{\rho}-0.85).
\end{equation}
Recommended production defaults: $\alpha=500$, $\delta=0$ (when quality filters are on), $\beta=1$, $\lambda=500$, $a_{\gamma}=0.5$, $b_{\gamma}=1$, $\kappa=10$.


\section{Scalability and implementation reference}
\label{sec:ch1-impl}

For $K_{\mathrm{samples}}=10^6$: multiprocessing, \texttt{float32} storage ($\sim$2.3\,GB), per-index RNG, in-place IPF, oversized candidate pool then filters, sampled evaluation. Candidate generation remains the throughput bottleneck; filters add a secondary pass over $K_{\mathrm{pool}}$.

\begin{table}[ht]
\centering
\caption{Chapter 1: key modules and outputs (\texttt{src/data/}).}
\label{tab:impl1}
\small
\begin{tabular}{@{}ll@{}}
\toprule
\textbf{Module / script} & \textbf{Role} \\
\midrule
\texttt{generate\_od\_dataset.py} & CLI orchestrator (prior, Gamma, filters) \\
\texttt{od\_analysis.py} & Base OD statistics \\
\texttt{od\_probability.py} & Distance prior $\varepsilon$ and $\mathbfp^{\mathrm{prior}}$ \\
\texttt{dirichlet\_sampler.py} & Dirichlet draw + demand reconstruction \\
\texttt{sparse\_mask.py} & Gamma heavy-tailed reweighting \\
\texttt{marginal\_perturbation.py} & Target marginals \\
\texttt{ipf.py} & IPF balancing \\
\texttt{od\_generator.py} & Parallel candidate generation \\
\texttt{quality\_filters.py} & Accept / reject pool filters \\
\texttt{od\_metrics.py} & Per-matrix and aggregate metrics \\
\texttt{dataset\_evaluation.py} & PCA, KMeans, diversity \\
\texttt{hyperparameter\_search.py} & Grid search over $(\alpha,\delta)$ \\
\midrule
\texttt{synthetic\_od.npy} & Output array shape $(K_{\mathrm{samples}},24,24)$ \\
\texttt{dataset\_metrics.json} & Evaluation summary \\
\bottomrule
\end{tabular}
\end{table}

\paragraph{Example command (1M samples).}
\begin{verbatim}
python generate_od_dataset.py \
  --input EstimatedODMatrix.npy \
  --samples 1000000 --oversample-factor 10 \
  --alpha 500 --perturbation 0 \
  --gamma-shape 0.5 --gamma-scale 1.0 \
  --epsilon-beta 1.0 --epsilon-lambda 500 \
  --memmap-save --skip-plots --skip-base-analysis
\end{verbatim}


%=============================================================================
\chapter{Turning-Count Observation and Dataset Construction}
\label{ch:turning}
%=============================================================================

\section{Network representation and turning movements}
\label{sec:ch2-network}

The road network is parsed from \texttt{sioux-falls.net.xml} (SUMO format) into a directed graph $\mathcal{G}=(\mathcal{V},\mathcal{E})$ using NetworkX. \textbf{No traffic simulation} is used---only static topology and connection rules: $|\mathcal{V}|=24$ junctions, $|\mathcal{E}|=76$ edges, 384 SUMO \texttt{<connection>} elements. Implementation: \texttt{src/data/network\_parser.py}.

A turning movement is a consecutive edge pair $(e^{\mathrm{in}},e^{\mathrm{out}})$ at junction $v$:
\begin{equation}
    \omega(e^{\mathrm{in}})=v=\alpha(e^{\mathrm{out}}), \qquad e^{\mathrm{in}}\neq e^{\mathrm{out}},
    \label{eq:turn_def}
\end{equation}
where $\omega(\cdot)$ and $\alpha(\cdot)$ denote head and tail junctions. Admissible movements yield $M=178$ indices stored in \texttt{turning\_id\_map} and exported as \texttt{turning\_movements.csv}. Implementation: \texttt{src/data/turning\_movements.py}.

\subsection{Binary vs.\ probabilistic route assignment}

Table~\ref{tab:route_gen_compare} compares the two forward models; \textbf{$K$-shortest-path + logit} (\texttt{generate\_turning\_counts.py}) is the production default.

\begin{table}[ht]
\centering
\caption{Comparison of route assignment and turning-count generation methods.}
\label{tab:route_gen_compare}
\small
\begin{tabular}{@{}p{3.2cm}p{4.8cm}p{4.8cm}@{}}
\toprule
\textbf{Criterion} & \textbf{Phase~1: binary shortest path} & \textbf{Phase~2: probabilistic (current)} \\
\midrule
Route set & Single shortest path $\pi_{ij}$ & Up to $K$ paths $\mathcal{R}_{ij}$ \\
Route choice & Deterministic & Multinomial logit on costs \\
$\Aturn$ entries & Binary $\{0,1\}$ & Fractional $[0,1]$ \\
$\mathrm{rank}(\Aturn)$ & 122 (nullity 454) & 174 (nullity 402) \\
Scalability at $10^6$ & Full in-memory load & Batched memory-mapped I/O \\
\bottomrule
\end{tabular}
\end{table}

Phase~1 understates route diversity and lowers $\mathrm{rank}(\Aturn)$. Phase~2 retains the linear forward model $\mathbfx=\Aturn\mathbfy$ while encoding route-choice structure via multinomial logit over a finite route catalog.


\section{Route enumeration and probabilistic assignment}
\label{sec:ch2-routes}

\subsection{OD pair indexing}

All $N^2=576$ zone pairs are indexed in row-major order:
\begin{equation}
    \ell(i,j)=(i-1)N+(j-1), \qquad y_{\ell(i,j)}=T_{ij}.
    \label{eq:od_index}
\end{equation}
Implementation: \texttt{src/data/od\_pairs.py}.

\subsection{Route catalog: $K$ shortest simple paths}

For each $(i,j)$, $i\neq j$, define $\mathcal{R}_{ij}=\{r_1,\ldots,r_{K_{ij}}\}$, $K_{ij}\le K$, from \texttt{networkx.shortest\_simple\_paths} in non-decreasing cost order. Route cost:
\begin{equation}
    c_r = \sum_{e\in r} w_e, \qquad w_e \in \{L_e, T_e\}.
    \label{eq:route_cost}
\end{equation}
Default $K=5$, $\theta=0.1$. Catalog saved as \texttt{route\_catalog.pkl}. Implementation: \texttt{src/data/k\_shortest\_paths.py}.

\begin{table}[ht]
\centering
\caption{Route enumeration and route choice: pros and cons.}
\label{tab:route_methods}
\small
\begin{tabular}{@{}p{3.2cm}p{4.2cm}p{4.2cm}@{}}
\toprule
\textbf{Method} & \textbf{Pros} & \textbf{Cons} \\
\midrule
$K$ shortest simple paths (chosen) & Cost-ordered; simulation-free; fixed $\Aturn$ & $K$ hyperparameter \\
Multinomial logit (chosen) & Closed form; one parameter $\theta$ & IIA; no congestion \\
User equilibrium & Realistic splits & Nonlinear; breaks fixed $\Aturn$ \\
Simulation-based & Congestion-aware & Slow at $10^6$ labels \\
\bottomrule
\end{tabular}
\end{table}

\subsection{Multinomial logit route choice}

Given costs $\{c_r\}$ on $\mathcal{R}_{ij}$:
\begin{align}
    U_r &= -\theta\, c_r, \label{eq:utility}\\
    P_r &= \frac{\exp(-\theta c_r)}{\sum_{r'\in\mathcal{R}_{ij}} \exp(-\theta c_{r'})},
    \label{eq:logit}
\end{align}
computed via stable softmax. Small $\theta$ increases route diversity and $\mathrm{rank}(\Aturn)$; large $\theta$ collapses toward Phase~1. Implementation: \texttt{src/data/route\_choice.py}.


\section{Fractional assignment matrix and forward model}
\label{sec:ch2-assignment}

The turn assignment matrix $\Aturn$ maps OD demand $\mathbfy\in\mathbb{R}_{\ge 0}^{N^2}$ to expected turning volumes $\mathbfx\in\mathbb{R}_{\ge 0}^{M}$. It is built \emph{once} from static network topology and route-choice probabilities, then reused for every synthetic OD replicate. This section gives the full construction procedure; the forward model that applies $\Aturn$ to $\mathbfy$ is stated at the end.


\subsection{Inputs, outputs, and indexing conventions}
\label{sec:ch2-aturn-io}

\paragraph{Inputs.}
\begin{itemize}[nosep]
    \item Parsed network $\mathcal{G}$ and turning index $\{1,\ldots,M\}$ with $M=178$ (\S\ref{sec:ch2-network}).
    \item OD column index $\ell(i,j)$ from \eqref{eq:od_index} for all $i\neq j$ (\S\ref{sec:ch2-routes}).
    \item Route catalog $\{\mathcal{R}_{ij}\}$ and logit probabilities $\{P_r\}_{r\in\mathcal{R}_{ij}}$ (\S\ref{sec:ch2-routes}).
\end{itemize}

\paragraph{Output.}
Matrix $\Aturn\in[0,1]^{M\times N^2}$, stored as \texttt{A\_turn.npy} with shape $(178,576)$ in \texttt{float32}. Row $m$ corresponds to turning movement $m$; column $\ell(i,j)$ corresponds to OD pair $(i,j)$. Self-loops ($i=j$) are omitted from routing and receive zero columns.

\paragraph{Matrix semantics.}
Entry $A_{\mathrm{turn},m,\ell(i,j)}$ is the \emph{expected fraction} of a unit trip from $i$ to $j$ that traverses turn $m$, after multinomial logit route choice over $\mathcal{R}_{ij}$. Columns are not constrained to sum to~1 because a route may use several turns; each column sum equals the expected number of distinct turns per trip for that OD pair.


\subsection{Turn incidence from a route path}
\label{sec:ch2-turn-incidence}

A route is an ordered edge list $r=(e_1,e_2,\ldots,e_L)$ with $L\ge 2$. At each internal step $k=1,\ldots,L-1$, consecutive edges $(e_k,e_{k+1})$ define a candidate turn at junction $\omega(e_k)=\alpha(e_{k+1})$. Let $\Theta_r\subset\{1,\ldots,M\}$ be the set of turn indices realized on $r$:
\begin{equation}
    \Theta_r = \bigl\{\, m : m = \mathrm{id}(e_k,e_{k+1}),\; k=1,\ldots,L-1 \,\bigr\},
    \label{eq:turn_set}
\end{equation}
where $\mathrm{id}(\cdot,\cdot)$ looks up $(e^{\mathrm{in}},e^{\mathrm{out}})$ in \texttt{turning\_id\_map}. Pairs not in the map (e.g.\ degenerate paths) are skipped. The \textbf{binary turn-incidence vector} is
\begin{equation}
    b_{r,m} = \mathbbm{1}[m\in\Theta_r], \qquad \mathbf{b}_r\in\{0,1\}^M.
    \label{eq:turn_vector}
\end{equation}
Implementation: \texttt{route\_turn\_indices} and \texttt{route\_turn\_vector} in \texttt{src/data/fractional\_assignment.py}. The Phase~1 binary matrix uses the same incidence rule on a single shortest path per OD (\texttt{src/data/build\_assignment\_matrix.py}).


\subsection{Column assembly for each OD pair}
\label{sec:ch2-column-assembly}

For each origin--destination pair $(i,j)$ with $i\neq j$, let $\ell=\ell(i,j)$. Initialize column $\mathbf{a}_\ell=\mathbf{0}\in\mathbb{R}^M$. For every route $r\in\mathcal{R}_{ij}$ with probability $P_r$ from \eqref{eq:logit},
\begin{equation}
    a_{\ell,m} \;\leftarrow\; a_{\ell,m} + P_r\, b_{r,m},
    \qquad m=1,\ldots,M.
    \label{eq:column_update}
\end{equation}
Equivalently, in vector form $\mathbf{a}_\ell = \sum_{r\in\mathcal{R}_{ij}} P_r\,\mathbf{b}_r$. Stacking columns yields the fractional assignment matrix:
\begin{equation}
    A_{\mathrm{turn},m,\ell} = \sum_{r\in\mathcal{R}_{ij}} P_r\, b_{r,m}.
    \label{eq:fractional_A}
\end{equation}
If $\mathcal{R}_{ij}=\emptyset$ (unreachable pair), column $\ell$ remains zero. With $K=5$ and $\theta=0.1$ on Sioux Falls, $\Aturn$ has 7{,}326 nonzeros (vs.\ 1{,}324 for the binary shortest-path matrix $\Aturn^{\mathrm{bin}}$).


\subsection{Binary baseline for comparison}
\label{sec:ch2-binary-baseline}

Phase~1 assigns one shortest path $\pi_{ij}$ per OD pair (\texttt{src/data/route\_assignment.py}) and sets
\begin{equation}
    A^{\mathrm{bin}}_{\mathrm{turn},m,\ell(i,j)} = \mathbbm{1}\bigl[m\in\Theta_{\pi_{ij}}\bigr],
    \label{eq:binary_A}
\end{equation}
so entries are in $\{0,1\}$. The production pipeline always builds both $\Aturn^{\mathrm{bin}}$ and $\Aturn$, compares ranks via \texttt{src/data/assignment\_rank.py}, and persists only the probabilistic $\Aturn$ for turning-count generation.


\subsection{Validation checks}
\label{sec:ch2-aturn-validation}

After construction, \texttt{validate\_fractional\_matrix} verifies:
\begin{itemize}[nosep]
    \item $0 \le A_{\mathrm{turn},m,\ell} \le 1$ for all $(m,\ell)$;
    \item column sums $s_\ell=\sum_m A_{\mathrm{turn},m,\ell}$ satisfy $s_\ell>0$ for routed OD pairs (typical mean $\bar{s}\approx 3.9$ turns per trip on Sioux Falls);
    \item probabilities within each $\mathcal{R}_{ij}$ sum to~1 (checked at route-choice stage).
\end{itemize}
Rank and condition number are reported in Table~\ref{tab:rank_compare}; $\mathrm{rank}(\Aturn)=174$ vs.\ $122$ for $\Aturn^{\mathrm{bin}}$.


\begin{figure}[ht]
\centering
\fbox{\begin{minipage}{0.94\linewidth}
\small
\textbf{Algorithm 2:} Construction of the turn assignment matrix $\Aturn$\\[0.4em]
\begin{enumerate}[leftmargin=*, nosep]
    \item \textbf{Parse network:} load \texttt{sioux-falls.net.xml}; build $\mathcal{G}$ and enumerate $M=178$ admissible turns from SUMO \texttt{<connection>} rules; export \texttt{turning\_movements.csv}
    \item \textbf{Index OD pairs:} map each $(i,j)$, $i\neq j$, to column $\ell(i,j)$; $N^2=576$ columns
    \item \textbf{Enumerate routes:} for each $(i,j)$, compute up to $K$ shortest simple paths $\mathcal{R}_{ij}$ in non-decreasing cost order; save \texttt{route\_catalog.pkl}
    \item \textbf{Route choice:} for each $\mathcal{R}_{ij}$, compute costs $c_r$, utilities $U_r=-\theta c_r$, and probabilities $P_r$ via stable softmax \eqref{eq:logit}
    \item \textbf{Turn incidence:} for each route $r=(e_1,\ldots,e_L)$, form $\Theta_r$ from consecutive edge pairs \eqref{eq:turn_set}; set $b_{r,m}=\mathbbm{1}[m\in\Theta_r]$
    \item \textbf{Assemble $\Aturn$:} for each $\ell=\ell(i,j)$, set $A_{\mathrm{turn},:, \ell}=\sum_{r\in\mathcal{R}_{ij}} P_r\,\mathbf{b}_r$ \eqref{eq:fractional_A}
    \item \textbf{Validate and compare:} check entries in $[0,1]$; build $\Aturn^{\mathrm{bin}}$ from shortest paths; report $\mathrm{rank}(\Aturn)$ and $\mathrm{rank}(\Aturn^{\mathrm{bin}})$
    \item \textbf{Persist:} save \texttt{A\_turn.npy} (shape $178\times 576$, \texttt{float32})
\end{enumerate}
\end{minipage}}
\caption{Turn assignment matrix generation (\texttt{generate\_turning\_counts.py}, stages 1--7; core logic in \texttt{fractional\_assignment.py}).}
\label{alg:aturn}
\end{figure}

Implementation: \texttt{src/data/fractional\_assignment.py} (\texttt{build\_fractional\_assignment\_matrix}), orchestrated by \texttt{generate\_turning\_counts.py}.


\subsection{Forward model}
\label{sec:ch2-forward}

Given a fixed $\Aturn$, the forward model for one replicate and for batches is:
\begin{equation}
    \mathbfx^{\mathrm{clean}} = \Aturn\,\mathbfy,
    \qquad
    \mathbfX^{\mathrm{clean}} = \mathbfY\,\Aturn^\top.
    \label{eq:forward_batch}
\end{equation}

\subsection{Identifiability and rank analysis}

\begin{table}[ht]
\centering
\caption{Assignment matrix rank comparison ($K=5$, $\theta=0.1$).}
\label{tab:rank_compare}
\begin{tabular}{@{}lcc@{}}
\toprule
\textbf{Quantity} & $\Aturn^{\mathrm{bin}}$ & $\Aturn$ (Phase~2) \\
\midrule
$\mathrm{rank}(\Aturn)$ & 122 & \textbf{174} \\
Nullity & 454 & 402 \\
$\kappa(\Aturn)$ & $\approx 14.8$ & $\approx 1.2\times 10^{12}$ \\
\bottomrule
\end{tabular}
\end{table}

Rank increases by 52, improving information content for Chapter~3. The inverse problem remains underdetermined ($\mathcal{N}=402$). Implementation: \texttt{src/data/assignment\_rank.py}.


\section{Observation noise, dataset assembly, and scalability}
\label{sec:ch2-dataset}

\subsection{Observation noise models}

Applied after clean synthesis. Real loop detectors and camera counts are noisy.

\paragraph{No noise.} $\mathbfx^{\mathrm{obs}} = \mathbfx^{\mathrm{clean}}$.

\paragraph{Gaussian noise.}
\begin{equation}
    x^{\mathrm{obs}}_m = x^{\mathrm{clean}}_m + \epsilon_m,
    \qquad \epsilon_m \sim \mathcal{N}\bigl(0,\sigma^2\,(x^{\mathrm{clean}}_m)_+\bigr),
    \label{eq:gauss_noise}
\end{equation}
with clipping at zero.

\paragraph{Poisson noise (default).}
\begin{equation}
    x^{\mathrm{obs}}_m \sim \mathrm{Poisson}\bigl(x^{\mathrm{clean}}_m\bigr),
    \label{eq:poisson_noise}
\end{equation}
with $\mathrm{Var}(x^{\mathrm{obs}}_m\mid\mathbfx^{\mathrm{clean}})=x^{\mathrm{clean}}_m$. Implementation: \texttt{src/data/observation\_noise.py}.

\begin{table}[ht]
\centering
\caption{Observation noise models: pros and cons.}
\label{tab:noise}
\small
\begin{tabular}{@{}p{3.5cm}p{4.5cm}p{4.5cm}@{}}
\toprule
\textbf{Model} & \textbf{Pros} & \textbf{Cons} \\
\midrule
None & Exact labels for sanity checks & Optimistic ML performance \\
Gaussian & Simple; additive & Negative counts without clipping \\
Poisson (chosen) & Standard for count data & Not for overdispersed sensors \\
Negative binomial & Handles overdispersion & Extra parameter \\
\bottomrule
\end{tabular}
\end{table}

\subsection{End-to-end pipeline and batched generation}

\begin{figure}[ht]
\centering
\fbox{\begin{minipage}{0.94\linewidth}
\small
\textbf{Network} $\to$ \textbf{$K$ paths} $\to$ \textbf{Logit} $\to$ \textbf{$\Aturn$ (Alg.~\ref{alg:aturn})}
$\to$ $\mathbfX^{\mathrm{clean}}=\mathbfY\Aturn^\top$ $\to$ \textbf{Poisson noise} $\to$ $\mathbfX^{\mathrm{obs}}$
\end{minipage}}
\caption{Chapter 2 turning-count pipeline (\texttt{generate\_turning\_counts.py}).}
\label{fig:phase2_pipeline}
\end{figure}

Loading all $10^6$ matrices into RAM caused OOM on a 15\,GB machine. The fix (\texttt{src/data/batch\_turning.py}) uses memory-mapped I/O, batched $\mathbfX=\mathbfY\Aturn^\top$, and streaming statistics. Benchmark: $10^6$ samples in $\approx$16\,s at $B=20{,}000$, peak RAM $\approx 60$\,MB per batch.

\begin{figure}[ht]
\centering
\fbox{\begin{minipage}{0.94\linewidth}
\small
\textbf{Algorithm 3:} Batched turning-count generation ($K_{\mathrm{samples}}=10^6$)\\[0.4em]
\begin{enumerate}[leftmargin=*, nosep]
    \item Run Algorithm~\ref{alg:aturn} once to obtain $\Aturn$
    \item Memory-map \texttt{synthetic\_od.npy}; never load full $\mathbfY$
    \item For each batch of size $B$: $\mathbfY^{(b)}\to\mathbfX^{\mathrm{clean},(b)}=\mathbfY^{(b)}\Aturn^\top$; apply Poisson; write rows
    \item Streaming mean/std; sampled pairwise diversity
\end{enumerate}
\end{minipage}}
\caption{Memory-safe batching (\texttt{src/data/batch\_turning.py}).}
\label{alg:batch_turning}
\end{figure}

\subsection{Validation, sensitivity, and evaluation highlights}

Forward consistency: $\|\mathbfx^{\mathrm{clean}}-\Aturn\mathbfy\|_\infty<10^{-3}$ per batch. Sensitivity grid over $(K,\theta)$ via \texttt{--run-sensitivity}. On $10^6$ synthetics: $\mathrm{rank}=174$; $\bar{d}_{\mathrm{OD}}\approx 545$; $\bar{d}_{\mathrm{turn}}\approx 894$; mean inter-turn correlation $\approx 0.01$.


\section{Implementation reference}
\label{sec:ch2-impl}

\begin{table}[ht]
\centering
\caption{Chapter 2: key modules (\texttt{src/data/}) and deliverables.}
\label{tab:impl2}
\small
\begin{tabular}{@{}ll@{}}
\toprule
\textbf{Module / file} & \textbf{Role} \\
\midrule
\texttt{generate\_turning\_counts.py} & CLI orchestrator \\
\texttt{k\_shortest\_paths.py} & $K$ shortest paths per OD \\
\texttt{route\_choice.py} & Multinomial logit \\
\texttt{build\_assignment\_matrix.py} & Binary baseline $\Aturn^{\mathrm{bin}}$ \\
\texttt{fractional\_assignment.py} & Fractional $\Aturn$ (Algorithm~\ref{alg:aturn}) \\
\texttt{batch\_turning.py} & Memory-mapped batched generation \\
\texttt{observation\_noise.py} & None / Gaussian / Poisson \\
\midrule
\texttt{A\_turn.npy} & $(178,576)$ fractional matrix \\
\texttt{turning\_counts.npy} & $(K_{\mathrm{samples}},178)$ clean counts \\
\texttt{turning\_counts\_noisy.npy} & Poisson observations \\
\texttt{route\_catalog.pkl} & Routes + logit probabilities \\
\bottomrule
\end{tabular}
\end{table}

\begin{verbatim}
python generate_turning_counts.py \
  --network sioux-falls.net.xml \
  --od outputs/od_generator/synthetic_od_1m.npy \
  --output-dir outputs/turning_counts_1m \
  --k_paths 5 --theta 0.1 --noise poisson \
  --batch-size 20000 --skip-plots
\end{verbatim}


%=============================================================================
\chapter{OD Estimation: Information Analysis, Models, and Benchmarking}
\label{ch:ml}
%=============================================================================

\section{Inverse problem and design principle}
\label{sec:ch3-inverse}

Given turning-count observations $\mathbfX\in\mathbb{R}^{K_{\mathrm{samples}}\times M}$, $M=178$, learn
\begin{equation}
    f:\mathbb{R}^{M}\to\mathbb{R}^{576}, \qquad \hat{\mathbfy}^{(k)} = f(\mathbfx^{\mathrm{obs},(k)}),
\end{equation}
approximating flattened OD $\mathbfy^{(k)}$. The forward model from Chapter~2 is
\begin{equation}
    \mathbfx^{\mathrm{clean}} = \Aturn\mathbfy, \qquad
    \mathbfx^{\mathrm{obs}} \sim \mathrm{Noise}(\mathbfx^{\mathrm{clean}}),
\end{equation}
with $\mathrm{rank}(\Aturn)=174$, nullity $\mathcal{N}=402$. The per-sample inverse map is underdetermined.

\begin{figure}[ht]
\centering
\fbox{\begin{minipage}{0.94\linewidth}
\small
$(\mathbfX,\mathbfY) \to$ statistics / PCA / MI / $\mathrm{rank}(\Aturn)$
$\to$ Ridge, PLS $\to$ RF, XGBoost $\to$ MLP, autoencoder $\to$ SHAP + leaderboard
\end{minipage}}
\caption{Chapter 3 pipeline (\texttt{train\_od\_estimator.py}). Classical models establish an information ceiling before deep models.}
\label{fig:phase3_pipeline}
\end{figure}

\paragraph{Design principle.}
Do \textbf{not} start with deep learning. First quantify recoverable OD information via $\mathrm{rank}(\Aturn)$, PCA dimensionality, and mutual information $\mathrm{MI}(x_m,y_\ell)$. If Ridge/PLS achieve low error, deep models test noise robustness---not representational capacity.

\paragraph{Underdetermined system vs.\ training set size.}
These are distinct: (i)~each sample has 178 observations for 576 unknowns ($\mathcal{N}=402$); (ii)~$K_{\mathrm{samples}}$ labeled pairs help \emph{learn} $f$, but do not remove per-sample ambiguity without regularization or priors.


\section{Dataset and information content analysis}
\label{sec:ch3-info}

Inputs $\mathbfX_{k,m}=x_m^{(k)}$ from \texttt{turning\_counts.npy}; targets $\mathbfY_{k,\ell}=y_\ell^{(k)}$ from flattened \texttt{synthetic\_od\_*.npy}. Implementation: \texttt{src/ml/dataset.py}.

Per-feature statistics (mean, std, CV), correlation matrices, and PCA ($n_{90\%}$ components for 90\% variance). Mutual information for sampled turn--OD pairs:
\begin{equation}
    \mathrm{MI}(x_m, y_\ell) = \iint p(x_m,y_\ell)\,\log\frac{p(x_m,y_\ell)}{p(x_m)p(y_\ell)}\,dx_m\,dy_\ell.
\end{equation}
Outputs: \texttt{important\_turns.csv}, \texttt{feature\_importance.csv}, \texttt{figures/pca\_X.png}, \texttt{singular\_value\_spectrum.png}. Implementation: \texttt{src/ml/information\_analysis.py}.


\section{Model zoo and selection rationale}
\label{sec:ch3-models}

\begin{table}[ht]
\centering
\caption{Model tiers: evaluation order, pros, and cons.}
\label{tab:model_order}
\footnotesize
\begin{tabular}{@{}p{2.2cm}p{2.8cm}p{3.8cm}p{3.8cm}@{}}
\toprule
\textbf{Tier} & \textbf{Models} & \textbf{Pros} & \textbf{Cons} \\
\midrule
Classical & Linear, Ridge, PLS & Fast; interpretable; PLS for $M<N^2$ & Linear; Lasso slow at 576 outputs \\
Tree & RF, ExtraTrees, XGBoost & Nonlinear; robust & Slow; 576-output regression \\
Neural & MLP, Residual MLP, AE & Flexible; latent $z\in\mathbb{R}^{16}$ & PyTorch; overfitting risk \\
\bottomrule
\end{tabular}
\end{table}

\paragraph{Why PLS.}
Partial Least Squares finds latent directions maximizing $\mathrm{cov}^2(\mathbfX\mathbfw_h,\mathbfY\mathbfq_h)$---well suited to collinear turning counts and multi-target regression.

\paragraph{Why Ridge before MLP.}
On clean Phase~2 data, $\mathbfx=\Aturn\mathbfy$ is exactly linear. Ridge with cross-validated $\alpha$ estimates the best regularized linear inverse for $\mathcal{N}=402$.

\paragraph{Autoencoder architecture.}
Optional two-stage: (i)~OD autoencoder $\mathbfy\to\mathbfz\in\mathbb{R}^{16}\to\hat{\mathbfy}$; (ii)~map $\mathbfx\to\mathbfz$, decode. Reduces output dimensionality when nullity is large.


\section{Physics constraints, loss functions, and evaluation metrics}
\label{sec:ch3-physics}

Non-negativity: $\hat{\mathbfy}\leftarrow\max(\hat{\mathbfy},\mathbf{0})$. Neural training may add production/attraction penalties:
\begin{equation}
    \mathcal{L}_{\mathrm{phys}} = \mathrm{MSE}(\hat{\mathbfy},\mathbfy)
    + \lambda_p \|\hat{\mathbfp}-\mathbfp\|_2^2
    + \lambda_a \|\hat{\mathbfa}-\mathbfa\|_2^2.
\end{equation}

\subsection{Loss functions}

\begin{table}[ht]
\centering
\caption{Loss functions for OD recovery.}
\label{tab:loss}
\begin{tabular}{@{}lll@{}}
\toprule
\textbf{Loss} & \textbf{Formula (single sample)} & \textbf{Use case} \\
\midrule
MSE & $\|\hat{\mathbfy}-\mathbfy\|_2^2$ & Default supervised \\
MAE & $\|\hat{\mathbfy}-\mathbfy\|_1$ & Robust to outliers \\
Forward-consistency & $\|\mathbfx-\Aturn\hat{\mathbfy}\|_2^2$ & Physics-informed \\
Combined & MSE $+\lambda\|\mathbfx-\Aturn\hat{\mathbfy}\|_2^2$ & Deployment \\
\bottomrule
\end{tabular}
\end{table}

Evaluation metrics on held-out test data:
\begin{align}
    \mathrm{MAE} &= \frac{1}{N^2}\sum_{\ell}|\hat{y}_\ell - y_\ell|, &
    \mathrm{RMSE} &= \sqrt{\frac{1}{N^2}\sum_{\ell}(\hat{y}_\ell-y_\ell)^2}, \\
    R^2 &= 1 - \frac{\sum_\ell(\hat{y}_\ell-y_\ell)^2}{\sum_\ell(y_\ell-\bar{y})^2}, &
    \rho_{\mathrm{OD}} &= \mathrm{corr}(\mathbfy,\hat{\mathbfy}), \\
    \|\mathbf{T}-\hat{\mathbf{T}}\|_F &= \sqrt{\sum_{i,j}(T_{ij}-\hat{T}_{ij})^2}. &
\end{align}
Production and attraction errors compare row and column sums. Implementation: \texttt{src/ml/metrics.py}.


\section{Training protocol, robustness, and explainability}
\label{sec:ch3-training}

80/10/10 train/validation/test split (seed-controlled). Standardize $\mathbfX$ on training statistics. Default \texttt{max\_samples=50000} for tractable benchmarking; Ridge/PLS can use larger $K_{\mathrm{samples}}$ (e.g.\ 100{,}000--1{,}000{,}000).

\paragraph{Robustness.}
Train on clean $\mathbfX^{\mathrm{clean}}$; evaluate on Poisson-noisy $\mathbfX^{\mathrm{obs}}$ at 0\%, 2\%, 5\%, 10\% levels.

\paragraph{SHAP.}
Tree-model SHAP values identify influential turning movements (\texttt{figures/shap\_summary.png}). Implementation: \texttt{src/ml/explainability.py}.

\paragraph{Training at scale.}
\begin{table}[ht]
\centering
\caption{Practical limits for $K_{\mathrm{samples}}=10^6$ on 15\,GB RAM.}
\label{tab:ml_scale}
\small
\begin{tabular}{@{}lll@{}}
\toprule
\textbf{Model} & \textbf{1M feasible?} & \textbf{Notes} \\
\midrule
Ridge, Linear, PLS & Yes & $\sim$3\,GB float32 arrays; minutes \\
Random Forest, XGBoost & Slow & Hours; high memory \\
Lasso, Elastic Net & No (default) & 576 separate sparse fits \\
MLP, Autoencoder & Slow & Use GPU; batch training \\
\bottomrule
\end{tabular}
\end{table}


\section{Benchmark leaderboard and deliverables}
\label{sec:ch3-benchmark}

\begin{table}[ht]
\centering
\caption{Chapter 3 deliverables (\texttt{outputs/ml/}).}
\label{tab:outputs3}
\begin{tabular}{@{}ll@{}}
\toprule
\textbf{File} & \textbf{Content} \\
\midrule
\texttt{best\_model.pkl} & Best model by test MAE \\
\texttt{benchmark\_results.csv} & Leaderboard (MAE, RMSE, $R^2$, marginals) \\
\texttt{feature\_importance.csv} & Ranked turning movements \\
\texttt{predictions.npy} & Test-set $\hat{\mathbfY}$ \\
\texttt{evaluation\_report.md} & Summary report \\
\bottomrule
\end{tabular}
\end{table}

\begin{table}[ht]
\centering
\caption{Chapter 3 modules (\texttt{src/ml/}).}
\label{tab:impl3}
\small
\begin{tabular}{@{}ll@{}}
\toprule
\textbf{Module} & \textbf{Role} \\
\midrule
\texttt{train\_od\_estimator.py} & CLI orchestrator \\
\texttt{dataset.py} & Load, split, standardize \\
\texttt{information\_analysis.py} & Stats, PCA, MI, rank \\
\texttt{models\_classical.py} & Linear, Ridge, PLS, trees \\
\texttt{models\_neural.py} & MLP, residual MLP, autoencoder \\
\texttt{benchmark.py} & Leaderboard and report \\
\bottomrule
\end{tabular}
\end{table}

\begin{verbatim}
# Full benchmark (use ODEstimation env Python)
/home/chetan/anaconda3/envs/ODEstimation/bin/python train_od_estimator.py \
  --turning outputs/turning_counts_1m/turning_counts.npy \
  --od outputs/od_generator/synthetic_od_1m.npy \
  --max-samples 50000

# Autoencoder on 100k samples
/home/chetan/anaconda3/envs/ODEstimation/bin/python train_od_estimator.py \
  --model autoencoder --max-samples 100000 --skip-analysis \
  --output-dir outputs/ml_autoencoder_100k
\end{verbatim}


\section{Future work and limitations}
\label{sec:ch3-future}

\begin{table}[ht]
\centering
\caption{Planned extensions beyond this document.}
\label{tab:future}
\begin{tabular}{@{}ll@{}}
\toprule
\textbf{Topic} & \textbf{Description} \\
\midrule
Real-data validation & Held-out counts / partial OD \\
Overdispersed noise & Negative binomial on $\mathbfx^{\mathrm{obs}}$ \\
User equilibrium & Nonlinear forward model \\
Uncertainty quantification & Bayesian / ensemble CIs on $\mathbfy$ \\
Deployment & Field turning-count inference pipeline \\
\bottomrule
\end{tabular}
\end{table}

\paragraph{Current limitations.}
\begin{itemize}[leftmargin=*]
    \item Single observed $\Tbase$; no temporal dynamics;
    \item Logit route choice without congestion feedback;
    \item Nullity $\mathcal{N}=402$ limits unique OD recovery without priors;
    \item Full 1M benchmark leaderboard not yet populated; subsampling used by default;
    \item Use \texttt{ODEstimation} conda env (not base \texttt{python}) for PyTorch/matplotlib compatibility.
\end{itemize}


%=============================================================================
\chapter*{Conclusion}
\addcontentsline{toc}{chapter}{Conclusion}

Chapter~1 produces diverse, structurally valid synthetic OD matrices via a distance-dependent gravity prior on full off-diagonal support, Dirichlet sampling, Gamma heavy-tailed reweighting, IPF balancing, and dataset quality filters---scalable to $10^6$ samples. Chapter~2 constructs turning-count observations through a route-choice-aware fractional operator $\Aturn$ ($\mathrm{rank}=174$, up from 122 under binary assignment), with optional Poisson noise and memory-mapped batch generation. Chapter~3 implements information analysis, classical baselines (Ridge, PLS), tree models, optional neural estimators, and a benchmark leaderboard via \texttt{src/ml/} and \texttt{train\_od\_estimator.py}. The three-chapter structure separates data generation from inverse estimation while maintaining consistent notation from $\Tbase$ through $\hat{\mathbfY}$.

\begin{thebibliography}{9}
\bibitem{furness1962}
R.~Furness, \emph{Trip forecasting}, Traffic Engineering and Control, 1962.

\bibitem{deming1940}
W.~E.~Deming and F.~F.~Stephan,
\emph{On a least squares adjustment of a sampled frequency table when the expected marginal totals are known},
Annals of Mathematical Statistics, 1940.

\bibitem{ortuzar2011}
J.~de~D.~Ort\'{u}zar and L.~G.~Willumsen,
\emph{Modelling Transport}, 4th ed., Wiley, 2011.

\bibitem{cascetta1984}
E.~Cascetta,
\emph{Estimation of trip matrices from traffic counts and survey data},
Transportation Research Part B, 1984.
\end{thebibliography}

\end{document}
